% =====================================================================
%  ALANYA — Dossier de conception
%  Étape 6 : identité et connexion par QR code
%
%  Compilation :  latexmk -pdf qr-code.tex
% =====================================================================
\documentclass[11pt,a4paper,oneside]{report}

\usepackage{alanya-conception}
\setAlanyaDocTitre{Identité et connexion QR}

\hypersetup{
  colorlinks=true,
  linkcolor=AlanyaIndigoStrong,
  urlcolor=AlanyaIndigoStrong,
  pdftitle={Alanya - Identite et connexion par QR code},
  pdfauthor={Equipe Alanya},
  pdfsubject={Conception - QR code d'identite et connexion multi-appareils}
}

\begin{document}

% =====================================================================
\begin{titlepage}
\thispagestyle{empty}
\begin{tikzpicture}[remember picture, overlay]
  \fill[AlanyaIndigo]
    (current page.north west) rectangle ([yshift=-6.4cm]current page.north east);
  \fill[AlanyaIndigoDark]
    ([yshift=-6.4cm]current page.north west) rectangle ([yshift=-6.0cm]current page.north east);
\end{tikzpicture}

\vspace*{1.1cm}
{\color{white}\sffamily
\begin{center}
  {\small\bfseries DOSSIER DE CONCEPTION}\\[6pt]
  {\footnotesize\color{AlanyaContainer} Document de travail --- à valider avant développement}
\end{center}}

\vspace{3.4cm}

\begin{center}\sffamily
  {\small\bfseries ÉTAPE 6}\\[10pt]
  {\color{AlanyaLine}\rule{0.6\linewidth}{0.6pt}}\\[20pt]
  {\fontsize{32}{36}\selectfont\bfseries\color{AlanyaIndigoDark}
   Identité et connexion par QR code}\\[10pt]
  {\large\color{AlanyaGray} Une seule brique de scan, deux usages}\\[6pt]
  {\normalsize\itshape\color{AlanyaGray}
   Ajout de contact instantané --- connexion multi-appareils}\\[22pt]
  {\color{AlanyaLine}\rule{0.6\linewidth}{0.6pt}}
\end{center}

\vspace{1.4cm}

\begin{center}
\begin{tikzpicture}[x=1cm, y=1cm]
  \foreach \x/\n in {0/Identité, 3.6/Connexion} {
    \draw[AlanyaIndigo, line width=1pt, rounded corners=1.5pt] (\x-0.22,-0.22) rectangle (\x+0.22,0.22);
    \fill[AlanyaIndigo] (\x-0.09,-0.09) rectangle (\x+0.09,0.09);
    \node[font=\sffamily\small, text=AlanyaGray, anchor=west] at (\x+0.42,0) {\n};
  }
\end{tikzpicture}
\end{center}

\vfill

\begin{center}\sffamily
  \begin{tabular}{@{}r@{\quad}l@{}}
    \thd{Périmètre}   & Talky (Flutter) \textbullet{} Alanya-Backend\\[3pt]
    \thd{Statut}      & Implémenté --- relu et corrigé\\[3pt]
    \thd{Dépend de}   & Rien --- indépendant des cinq premiers volets\\[3pt]
    \thd{Migration}   & \code{026\_qr\_appareils.sql}\\[3pt]
    \thd{Date}        & \today\\
  \end{tabular}
\end{center}

\vspace{0.8cm}
\begin{center}\color{AlanyaGray}\small\sffamily
  Document confidentiel --- diffusion restreinte
\end{center}
\end{titlepage}

% =====================================================================
\chapter*{Avant-propos}
\addcontentsline{toc}{chapter}{Avant-propos}
\markboth{Avant-propos}{}
\thispagestyle{plain}

Deux fonctionnalités demandées, toutes deux fondées sur le scan d'un QR code~: un code
d'identité personnel qui ajoute instantanément un contact préféré, et une connexion
multi-appareils façon « appareil de confiance qui en autorise un nouveau ». Elles sont
traitées dans un seul dossier parce qu'elles partagent une vraie infrastructure --- le
composant scanner, le générateur de QR stylé, l'enveloppe de résolution d'un code scanné ---
au même titre que compte officiel et diffusion l'ont été à l'étape~3.

\begin{note}[Comment lire ce document]
La \textbf{partie I} décrit la fonctionnalité telle que l'utilisateur et le client la
voient. La \textbf{partie II} décrit comment on la construit. On peut valider la partie~I
sans lire la partie~II.
\end{note}

\tableofcontents

% =====================================================================
\part{Approche naïve}
% =====================================================================

% Fragment partage : inclus par qr-code.tex (dossier autonome)
% et par dossier-conception-alanya.tex (dossier client assemble).
% Ne pas y mettre de preambule ni de \part : c'est l'appelant qui les pose.

\chapter{Le besoin}

\section{Ce qui existe déjà}

Alanya n'a aujourd'hui \textbf{qu'un seul chemin} pour ajouter un contact préféré~: la
recherche manuelle par nom, pseudo ou Alanya ID, depuis la feuille d'ajout du profil.
Aucun lien d'invitation, aucun code, aucun geste plus rapide.

\begin{constat}[Vérifié dans le code]
\code{lib/widgets/add\_contact\_sheet.dart} interroge d'abord le cache local
(\code{filterUsersBySearch}), puis \code{GET~/users/search} après 400~ms sans résultat.
L'ajout appelle \code{POST~/contacts/:userId} --- la même route que celle que ce volet
réutilisera pour l'ajout par QR.
\end{constat}

Côté connexion, Alanya n'a \textbf{qu'un seul appareil implicite}~: se connecter sur un
second téléphone fonctionne déjà techniquement (rien ne l'empêche), mais rien ne le montre
non plus. Il n'existe aucun écran « mes appareils », aucun moyen de savoir où son compte est
ouvert, ni de fermer une session à distance.

\begin{constat}[Vérifié dans le code]
\code{lib/screens/profile/account\_security\_screen.dart} ne contient que deux
sections~: \emph{Email} et \emph{Changer le mot de passe}. Rien sur les appareils.
\end{constat}

\section{Ce qu'on ajoute}

Deux fonctionnalités, toutes deux fondées sur un QR code~:

\begin{enumerate}[leftmargin=1.6em]
  \item \textbf{Un code personnel éphémère.} Chaque utilisateur génère, en ouvrant
        « Mon code », un QR à son identité --- logo Alanya au centre, traits arrondis,
        une vraie carte de visite plutôt qu'un carré noir et blanc. Le scanner d'un autre
        utilisateur reconnaît la personne et l'ajoute \emph{instantanément} en contact
        préféré. Le code vit \textbf{dix minutes} et s'éteint au premier scan réussi~;
        l'écran en régénère un neuf tout seul, et le propriétaire est prévenu --- et
        invité à ajouter l'autre \textbf{en retour} --- chaque fois que son code est
        utilisé.
  \item \textbf{La connexion d'un second appareil.} Un nouvel appareil affiche un code ; un
        appareil déjà connecté le scanne pour l'autoriser. L'utilisateur peut ensuite voir,
        depuis son compte, la liste de tous les appareils connectés --- et en déconnecter un
        à distance.
\end{enumerate}

\begin{note}[Une seule brique, deux usages]
Les deux fonctionnalités ne partagent pas qu'un mot dans leur nom. Elles partagent un
\textbf{composant scanner}, un \textbf{générateur de QR stylé}, et une manière commune de
résoudre ce qu'un code scanné signifie. C'est pour cette raison qu'elles sont conçues dans
un seul volet --- le même raisonnement qui a réuni compte officiel et diffusion dans
l'étape~3.
\end{note}

\section{Les règles produit}

\begin{center}
\begin{tabularx}{\linewidth}{@{}p{4.6cm}X@{}}
\toprule
\thd{Qui peut scanner} & Un utilisateur \textbf{déjà connecté}, dans les deux cas. Scanner
n'est jamais un geste public ; c'est toujours un compte qui en reconnaît un autre. \\
\thd{Vie du code personnel} & \textbf{Dix minutes, une seule personne.} Généré à la
demande à l'ouverture de « Mon code », consommé au premier scan réussi. L'écran affiche
un compte à rebours et régénère automatiquement --- à l'expiration comme après
consommation. \\
\thd{Ajout de contact} & \textbf{Instantané}, sans écran de confirmation. Le filet de
sécurité est une carte de résultat annulable, pas une étape supplémentaire. \\
\thd{Retour au propriétaire} & Prévenu à \textbf{chaque scan} de son code, où qu'il soit
dans l'application, et invité à ajouter l'autre en retour --- notification si l'app est
en arrière-plan, simple information si le lien existait déjà. \\
\thd{Contacts ajoutés par QR} & Reconnaissables partout~: pastille sur l'avatar, mention
datée sur la fiche, filtre « Par QR » dans les contacts préférés \emph{et} dans la liste
des discussions. \\
\thd{Note de contexte} & Optionnelle, saisie juste après le scan --- « rencontré au salon
de Douala » --- ou au moment d'ajouter en retour, et relue sur la fiche du contact. Elle
appartient à la relation~: l'autre ne la voit jamais. \\
\thd{Codes permanents} & \textbf{Réservés aux futurs comptes business} (volet~1). Aucun
compte personnel n'a de code qui dure. \\
\thd{Connexion d'un appareil} & \textbf{Jamais instantanée.} L'appareil qui scanne voit
toujours l'identité de l'appareil qu'il autorise, et doit confirmer explicitement. \\
\thd{Durée du QR de connexion} & Environ \textbf{90~secondes}, puis régénération
automatique --- comme un code à usage unique, pas comme une carte de visite. \\
\thd{Emplacements} & Profil (mon code), feuille d'ajout de contact (scanner), Compte et
sécurité (appareils connectés), écran de connexion (nouvel appareil). \\
\bottomrule
\end{tabularx}
\end{center}

\begin{note}[Pourquoi un code qui meurt]
Un code permanent est une affiche~; un code de dix minutes est une poignée de main. Le
code personnel sert à des rencontres --- on le montre, on l'envoie, la personne en face
scanne, c'est fini. L'usage unique garantit que le code n'ajoute que la personne à qui on
vient de le tendre~: une capture d'écran qui circule ensuite ne mène plus nulle part. Et
comme l'écran régénère aussitôt, montrer son code à trois personnes d'affilée reste
trois gestes de \emph{leur} côté, zéro du sien.
\end{note}

\begin{constat}[Vérifié dans le code --- pourquoi le permanent est verrouillé]
L'infrastructure du code permanent (\code{users.qr\_public\_id}) existe et reste en
place, mais son garde refuse aujourd'hui tout le monde --- et pour une raison précise~:
la colonne \code{users.type\_compte}, qu'on prendrait volontiers pour un type de compte
métier, est en réalité un \textbf{rôle d'administration} (0~utilisateur, 1~admin,
2~super-admin, cf. la route \code{PUT~/users/:id/role} de \code{src/routes/admin.js},
côté Alanya-Backend). Rien dans la base ne sait encore dire « ce compte est un
commerce ». Le volet~1 (socle de compte) apportera \code{account\_type}~; le verrou des
codes permanents s'y branchera alors.
\end{constat}

\chapter{Mon code personnel}

\section{Vue d'ensemble}

Un seul écran, deux volets accessibles d'un geste~: \emph{Mon code} et \emph{Scanner} ---
comme on bascule entre appareil photo avant et arrière.

\begin{center}
  \imageOuCadre{maquettes/qr-identite.png}{\linewidth}%
    {Mon code (médaillon, logo au centre, partage et régénération) et l'écran Scanner
     (reconnaissance instantanée, succès doux et annulable)\\
     \texttt{maquettes/qr-identite.png}}
\end{center}

\begin{note}[Maquette interactive]
\href{https://claude.ai/code/artifact/69e9cbec-ff13-48fc-903f-673648874b2a}%
{\texttt{claude.ai/code/artifact/69e9cbec}}\\
Source versionnée~: \code{docs/architecture/maquettes/qr-identite.html}
\end{note}

\section{Écran --- Mon code}

Une carte blanche~: médaillon d'avatar débordant le haut, nom, Alanya ID, puis le QR
code lui-même --- fond clair, points arrondis à la couleur de marque, logo Alanya au
centre. Sous le code, un \textbf{compte à rebours} (« Expire dans 9:41 ») et la mention
« Valable 10~minutes et pour une seule personne ». En dessous, deux actions~:
\textbf{Partager} et \textbf{Nouveau code}.

À zéro --- comme après un scan réussi --- l'écran régénère un code neuf de lui-même~:
l'utilisateur n'a jamais un code mort sous les yeux, exactement le comportement déjà
retenu pour le QR de connexion (chapitre~\ref{ch:qr-naif-connexion}). « Nouveau code »
reste disponible pour qui veut couper court sans attendre~; aucun dialogue de
confirmation, renouveler un code qui expire tout seul n'a rien d'irréversible.

\begin{note}[Pourquoi le code reste toujours clair]
Même en thème sombre, la carte du QR reste sur fond blanc. Un code sur fond sombre scanne
moins bien --- la fiabilité du scan prime sur la cohérence du thème, à cet endroit
précisément.
\end{note}

\section{Être prévenu, ajouter en retour}
\label{sec:qr-naif-retour}

Quand quelqu'un utilise votre code --- scan ou lien ---, vous êtes prévenu
immédiatement~: un dialogue propose d'ajouter la personne \textbf{en retour} (« untel
vous a ajouté à ses contacts préférés avec votre code QR. L'ajouter en retour ? »,
\emph{Ajouter} / \emph{Non merci}). Si l'application est en arrière-plan ou fermée, une
notification prend le relais. Et si vous aviez déjà cette personne en contact, pas de
question~: une simple information.

\begin{note}[Le dialogue n'habite aucun écran]
Le code a pu être partagé par lien~: à l'instant où il est utilisé, son propriétaire
peut être n'importe où dans l'application --- pas forcément sur l'écran « Mon code ». Le
dialogue est donc global, affiché où que l'on soit.
\end{note}

L'écran « Mon code », lui, réagit à sa manière~: le code venant d'être consommé, il en
affiche un neuf aussitôt --- le propriétaire peut enchaîner plusieurs personnes sans un
geste.

\section{Reconnaître un contact ajouté par QR}

Un contact ajouté par QR reste reconnaissable après coup, partout où il apparaît~:

\begin{enumerate}[leftmargin=1.6em]
  \item une \textbf{pastille QR} au coin bas-gauche de l'avatar --- le bas-droit reste
        réservé au point de présence en ligne, les deux cohabitent --- dans
        \textbf{toutes} les listes~: contacts préférés, feuille d'ajout, nouvelle
        discussion, sélection de membres, partage de contact\ldots{} et jusqu'à la liste
        des discussions~;
  \item une \textbf{mention datée} sur la fiche du contact~: « Ajouté par QR code ·
        30~juil.~2026 »~;
  \item un \textbf{filtre} « Tous / Par QR » sur l'écran des contacts préférés, chaque
        option affichant son compte.
\end{enumerate}

L'ajout \emph{en retour} est marqué de la même façon~: les deux directions du lien
portent la même origine.

\begin{note}[Maquette interactive]
\href{https://claude.ai/code/artifact/0e095614-eb7c-43bd-8f14-b0d15e61c139}%
{\texttt{claude.ai/code/artifact/0e095614}}\\
Source versionnée~: \code{docs/architecture/maquettes/qr-distinction.html}
\end{note}


La liste des discussions porte le même filtre « Par QR » que l'écran des contacts
préférés~: un chip parmi les filtres existants (Tous, Discussions, Groupes…), qui ne
retient que les conversations avec un contact rencontré par QR.

Enfin, la carte de résultat du scan propose un champ facultatif~: une \textbf{note de
contexte}, quelques mots pour se souvenir de la rencontre. Elle se retrouve sur la fiche
du contact, en italique sous la mention d'origine --- et nulle part ailleurs~: c'est un
aide-mémoire personnel, jamais montré à l'intéressé.

\section{Partager son code}

\textbf{Partager} ouvre un choix explicite~: le \textbf{lien} (cliquable, accompagné de
l'Alanya ID) ou l'\textbf{image} (la carte telle qu'affichée, compte à rebours compris).
Dans les deux cas la légende annonce la couleur --- « Ce code est valable 10~minutes et
pour une seule personne » --- suivie de l'Alanya ID~: le code meurt en dix minutes,
l'Alanya ID reste le chemin permanent, les deux voyagent ensemble.

\begin{note}[Pas d'enregistrement dans la galerie]
Enregistrer la carte en image n'est volontairement pas proposé~: une image morte au bout
de dix minutes finirait en code cassé au fond d'une pellicule. L'option reviendra avec
les codes permanents des comptes business --- eux supportent l'impression et
l'affichage.
\end{note}

\section{Écran --- Scanner}

Plein écran, viseur au centre, façon appareil photo. Dès qu'un code d'identité est
reconnu, une \textbf{carte de résultat} se pose en bas de l'écran~: la personne ajoutée
(avatar, nom), « Ajouté à vos contacts », et trois suites possibles ---
\textbf{Message}, \textbf{Voir détails}, \textbf{Annuler}. La caméra reste vivante
derrière la carte~: on peut enchaîner plusieurs personnes sans ressortir de l'écran. Pas
d'écran intermédiaire, pas de confirmation à taper.

Trois cas particuliers, traités comme des succès doux plutôt que des erreurs~: scanner
son propre code (rien ne se passe, un message discret l'indique --- et le code
\emph{survit}), scanner un contact déjà favori (la carte s'affiche quand même, sans
bouton Annuler~: rien n'a été écrit), scanner un code expiré, déjà utilisé ou illisible
(message clair, la caméra reste active).

Un bouton \textbf{Importer une image} couvre enfin le code reçu en capture d'écran ou en
photo~: l'image choisie dans la galerie est analysée exactement comme le cadre de la
caméra --- même résultat, même carte.

\section{Points d'entrée}

\begin{enumerate}[leftmargin=1.6em]
  \item Le profil, une icône QR à côté de l'Alanya ID
        (\code{lib/screens/profile/profile\_screen.dart}, en-tête).
  \item La feuille d'ajout de contact, un bouton « Scanner un code » au-dessus de la
        recherche existante (\code{lib/widgets/add\_contact\_sheet.dart}).
\end{enumerate}

\chapter{Connexion d'un second appareil}
\label{ch:qr-naif-connexion}

\section{Le principe}

Contrairement à l'ajout de contact, une connexion ne s'approuve \textbf{jamais} au simple
scan. Le nouvel appareil affiche un code ; c'est l'appareil \emph{déjà} connecté qui le
scanne et confirme --- jamais l'inverse.

\begin{alerte}[Pourquoi une confirmation, jamais une reconnaissance instantanée]
Si le scan suffisait à connecter, un attaquant pourrait afficher le QR de connexion d'un
appareil qu'il contrôle --- sur un écran, un site, une affiche --- et attendre qu'une
victime le scanne par réflexe avec son propre compte déjà connecté. C'est une attaque
documentée contre les connexions par QR (le «~quishing~»). L'écran de confirmation, qui
montre \emph{quel} appareil demande la connexion avant de l'accepter, est le seul rempart~:
il doit toujours s'afficher, sans exception ni raccourci.
\end{alerte}

\section{Écran --- Se connecter par QR}

Sur le \textbf{nouvel} appareil, non encore connecté~: un code, un compte à rebours, un
statut (« En attente de scan… »). Le code se régénère seul à l'expiration --- l'utilisateur
n'a jamais un code mort sous les yeux.

\begin{center}
  \imageOuCadre{maquettes/qr-connexion-parcours.png}{\linewidth}%
    {La poignée de main~: le nouvel appareil affiche, l'appareil de confiance confirme\\
     \texttt{maquettes/qr-connexion-parcours.png}}
\end{center}

\begin{note}[Maquette interactive]
\href{https://claude.ai/code/artifact/a69b2050-b2dc-4aa1-a204-59ada2526a86}%
{\texttt{claude.ai/code/artifact/a69b2050}}\\
Source versionnée~: \code{docs/architecture/maquettes/qr-connexion.html}
\end{note}

\section{Écran --- Confirmer la connexion}

Sur l'appareil \textbf{déjà} connecté, après le scan~: une fiche décrivant ce qui vient
d'être scanné --- type d'appareil, ville approximative, à l'instant --- suivie de deux
boutons, \textbf{Refuser} et \textbf{Confirmer}. Un avertissement rappelle de refuser et de
changer son mot de passe en cas de doute.

\begin{note}[Pourquoi l'adresse compte plus que le nom de l'appareil]
Le nom et la plateforme sont \emph{annoncés} par l'appareil qui demande la connexion~: rien
n'empêche un attaquant d'appeler le sien « Alanya --- Vérification de sécurité ». La seule
information que le serveur \emph{constate} lui-même est l'adresse d'où part la demande.
C'est donc la seule sur laquelle fonder sa décision, et l'écran le dit explicitement.
\end{note}

\begin{note}[Une ville, pas une adresse technique]
Cette adresse est traduite en lieu approximatif --- « Douala, Cameroun » --- avant d'être
montrée. Sous sa forme brute (\code{41.202.219.14}), elle n'apprendrait rien à qui n'est pas
technicien~: on afficherait la seule information fiable de l'écran sous la seule forme où
elle est inutilisable. Sous forme de ville, le doute devient une évidence --- on sait
immédiatement si la demande part d'où l'on se trouve. Si la traduction n'aboutit pas,
l'adresse brute reste affichée plutôt que rien.
\end{note}

\section{Écran --- Mes appareils connectés}

Un écran à part, indépendant de toute connexion en cours, atteint depuis \emph{Compte et
sécurité}~: tous les appareils sur lesquels le compte est ouvert --- pas seulement ceux
ajoutés par QR --- avec un repère « Cet appareil » et un bouton \textbf{Déconnecter} par
ligne, à effet immédiat.

\begin{center}
  \imageOuCadre{maquettes/qr-connexion-appareils.png}{\linewidth}%
    {Mes appareils connectés~: tous les logins, pas seulement ceux par QR\\
     \texttt{maquettes/qr-connexion-appareils.png}}
\end{center}

\begin{note}[Maquette interactive]
\href{https://claude.ai/code/artifact/a69b2050-b2dc-4aa1-a204-59ada2526a86}%
{\texttt{claude.ai/code/artifact/a69b2050}} --- second panneau de la même maquette\\
Source versionnée~: \code{docs/architecture/maquettes/qr-connexion.html}
\end{note}

\section{Points d'entrée}

\begin{enumerate}[leftmargin=1.6em]
  \item L'écran de connexion, un bouton « Se connecter avec un code QR »
        (\code{screens/authentification/\allowbreak login\_screen.dart}), pour l'appareil
        qui n'est pas encore connecté.
  \item \emph{Compte et sécurité}, une nouvelle section « Appareils connectés »
        (\code{screens/profile/\allowbreak account\_security\_\allowbreak screen.dart}),
        pour lier un nouvel appareil ou consulter la liste depuis un appareil déjà
        connecté.
\end{enumerate}


% =====================================================================
\part{Approche technique}
% =====================================================================

% Fragment partage : inclus par qr-code.tex (dossier autonome)
% et par dossier-conception-alanya.tex (dossier client assemble).
% Ne pas y mettre de preambule ni de \part : c'est l'appelant qui les pose.

\chapter{Enveloppe de payload et résolution}
\label{ch:qr-payload}

Les deux fonctionnalités encodent la même chose dans leur QR~: une URL portant un
\textbf{jeton opaque}, jamais une donnée brute. Ni l'\code{alanyaID} interne
(auto-incrémenté, donc énumérable), ni l'Alanya ID (\code{alanyaPhone}) --- ce dernier est
aussi un identifiant de connexion, et l'exposer sur un support scannable par n'importe qui
affaiblirait sa fonction actuelle.

\begin{center}
\begin{tabularx}{\linewidth}{@{}p{5.4cm}X@{}}
\toprule
\thd{Chemin encodé} & \thd{Ce qu'il résout} \\
\midrule
\code{/q/u/:qrPublicId} & un utilisateur --- ajout de contact
(\code{users.qr\_public\_id}) \\
\code{/q/l/:sessionId.:scanSecret} & une session de connexion en attente, en mémoire
(chapitre~\ref{ch:qr-ephemere}) \\
\bottomrule
\end{tabularx}
\end{center}

\begin{note}[Une seule base d'URL, dans un seul fichier]
\code{src/utils/qrToken.js} porte \code{QR\_BASE\_URL} (variable d'environnement, par
défaut \code{https://\allowbreak www.\allowbreak alanya237.com}) ainsi que les deux
fabricants de payload, \code{identityPayload} et \code{loginPayload}. Ils vivent ensemble
et non dans leurs contrôleurs respectifs parce qu'un seul contrôleur, \code{qrController},
\emph{reparse} les deux~: deux constantes séparées ont déjà divergé une fois, laissant les
QR de connexion sur l'ancien domaine. Le même fichier porte aussi
\code{generateOpaqueToken} et \code{secretMatches}, la comparaison à temps constant des
secrets. Le point qui sépare \code{sessionId} de \code{scanSecret} n'est jamais ambigu~:
les deux jetons sont en \code{base64url}, dont l'alphabet ne contient pas le point.
\end{note}

Les deux types partagent un seul point d'entrée~:

\begin{center}
\begin{tabularx}{\linewidth}{@{}lX@{}}
\toprule
\mPOST{}~\code{/api/qr/resolve} & Corps~: \code{\{ payload \}}, la chaîne \textbf{brute}
scannée --- le client ne l'interprète jamais. Réponse~: \code{\{type:'contact',
\ldots{}\}} (chapitre~\ref{ch:qr-contact}) ou \code{\{type:'login', session\}}. \\
\bottomrule
\end{tabularx}
\end{center}

\begin{note}[Le domaine n'est pas vérifié au parsing]
\code{\_parsePayload} n'extrait que le \emph{chemin} de l'URL scannée~: l'autorité de
résolution est le jeton, pas l'hôte. Un QR déjà imprimé sous un ancien domaine continue
donc de fonctionner, et un QR forgé sous un autre domaine ne gagne rien --- il faudrait
qu'il porte un jeton valide, ce que le contrôle d'autorisation vérifie ensuite.
\end{note}

\begin{alerte}[Le scanneur doit toujours être authentifié]
Les deux flux exigent que \textbf{qui scanne} soit déjà connecté --- pas seulement qui
possède un compte. Scanner un contact sans être soi-même connecté n'a pas de sens (à qui
l'ajouter~?) ; scanner une session de connexion sans l'être serait la faille elle-même. Le
middleware \code{authCustom} garde donc \code{/api/qr/resolve} au même titre que le reste
de l'API, doublé de \code{qrResolveLimiter} (30~résolutions par minute et par~IP) --- une
défense en profondeur contre l'énumération des jetons, pas une contrainte fonctionnelle.
\end{alerte}

\chapter{Modèle de données}
\label{ch:qr-donnees}

\section{Le jeton d'identité}

Une colonne sur \code{users}, indépendante de \code{alanyaPhone}~: régénérable sans effet
sur la connexion, sans capacité de connexion elle-même.

\section{La nature de \code{device\_id}}
\label{sec:qr-device-id}

\begin{note}[Un identifiant matériel, pas un UUID d'application]
\code{lib/\allowbreak core/\allowbreak services/\allowbreak storage\_service.dart}, méthode
\code{getOrCreateDeviceId()}, génère un UUID
aléatoire persisté par \code{flutter\_secure\_storage} --- stable tant que l'app reste
installée, mais \textbf{réinitialisé à la réinstallation}. Un appareil réinstallé
apparaîtrait alors comme un second appareil fantôme dans la liste. Le \code{device\_id} de
la table \code{appareils} est donc un identifiant \textbf{matériel}, calculé par
\code{TalkyApiClient.\_currentHardwareId()} et distinct de l'UUID applicatif, que le push
et le socket continuent d'utiliser de leur côté.
\end{note}

\begin{center}
\begin{tabularx}{\linewidth}{@{}p{2.5cm}X@{}}
\toprule
\thd{Plateforme} & \thd{Source de \code{device\_id}} \\
\midrule
Android & \code{Settings.\allowbreak Secure.\allowbreak ANDROID\_ID}, lu par un
\textbf{canal natif Kotlin} --- \code{MethodChannel} \code{com.alanya/\allowbreak
device\_identity}, méthode \code{getAndroidId}, déclaré dans
\code{android/\allowbreak app/\allowbreak src/main/\allowbreak kotlin/\allowbreak
com/\allowbreak alanya237/\allowbreak alanya/\allowbreak MainActivity.kt}. Unique par
(appareil, signature d'app, utilisateur) et survivant à une réinstallation. \\
iOS & \code{identifierForVendor}, via \code{device\_info\_plus}. \\
Autres & repli direct sur l'UUID applicatif. \\
\bottomrule
\end{tabularx}
\end{center}

\begin{alerte}[Pourquoi un canal natif et non \code{device\_info\_plus}]
\code{device\_info\_plus} n'expose plus \code{ANDROID\_ID} depuis sa \textbf{v9} --- le
paquet reste une dépendance et sert toujours au libellé lisible de l'appareil
(\code{\_currentDeviceModel()}) et à l'\code{identifierForVendor} d'iOS, mais il ne peut
plus fournir la clé d'appareil sur Android. Le piège à éviter est
\code{AndroidDeviceInfo.id}, qu'on prendrait volontiers pour son remplaçant~: ce champ est
\code{Build.ID}, c'est-à-dire l'\textbf{empreinte du firmware}~: identique sur tous les
téléphones d'un même modèle, et modifiée à chaque mise à jour système. Comme l'unicité
porte sur le couple (\code{alanyaID}, \code{device\_id}), deux téléphones du même modèle
appartenant au \emph{même} compte se retrouveraient sur une seule ligne --- en déconnecter
un déconnecterait l'autre --- pendant qu'une simple mise à jour ferait réapparaître chaque
téléphone comme neuf. Et un identifiant que tout possesseur du même modèle connaît n'est
plus un identifiant. Strictement inutilisable comme clé d'appareil~; d'où le canal natif.
\end{alerte}

\begin{note}[La chaîne de repli, jusqu'au bout]
Si \code{ANDROID\_ID} est indisponible (exception ou valeur vide), et hors Android/iOS
(web, macOS, Windows, Linux), \code{\_currentHardwareId()} retombe sur l'UUID applicatif
\code{StorageService().getOrCreateDeviceId()} --- puis, en dernier recours, sur
\code{'INDEFINI'}. Mieux vaut un identifiant par installation qu'aucun~: sans
\code{device\_id}, aucune ligne \code{appareils} n'est créée, donc aucun jeton n'est émis
(chapitre~\ref{ch:revocation}) et la connexion échoue. Le prix de ce repli est la seule
imperfection assumée du dispositif~: sur ces plateformes, une réinstallation crée une
seconde ligne.
\end{note}

\begin{alerte}[Asymétrie de persistance entre plateformes]
Le Keychain iOS --- et donc potentiellement l'UUID applicatif lui-même --- survit
généralement à une réinstallation ; les données d'application Android, dont l'UUID
applicatif, sont effacées à la désinstallation. \code{ANDROID\_ID} est donc le seul des
deux identifiants qui apporte une vraie amélioration côté Android. Côté iOS,
\code{identifierForVendor} ne survit pas à la désinstallation de la \emph{dernière}
application du vendeur~: le téléphone réapparaît alors comme un nouvel appareil. Limite
assumée, et signalée à l'utilisateur dans l'écran « Appareils connectés ».
\end{alerte}

Conséquence assumée~: \code{appareils.device\_id} et \code{user\_push\_devices.deviceId}
ne sont \textbf{pas} directement joignables --- ce sont deux identifiants différents. La
table \code{appareils} porte donc un \code{push\_device\_id} nullable, réservé à
l'affichage optionnel du statut « notifications activées » dans la liste des appareils.
\textbf{Aucun code ne l'écrit à ce jour}~: la colonne est posée pour cet usage futur, et
l'écran doit se comporter comme si elle était toujours vide.

\section{Le jeton et la table \code{appareils}}

\begin{lstlisting}[language=SQL, caption={\code{migrations/026\_qr\_appareils.sql}}]
-- Identite QR (users.qr_public_id) + registre des appareils connectes
-- (appareils), alimente par toute connexion (login, register, QR).
--
-- MySQL 8 ne supporte pas ADD COLUMN IF NOT EXISTS (cf. migration 008) : si
-- relancee apres un premier passage reussi, ignorer les erreurs 1060
-- (Duplicate column name) / 1061 (Duplicate key name).

ALTER TABLE users
  ADD COLUMN qr_public_id VARCHAR(64) NULL,
  ADD UNIQUE KEY uq_users_qr_public_id (qr_public_id);

CREATE TABLE IF NOT EXISTS appareils (
  id BIGINT AUTO_INCREMENT PRIMARY KEY,
  alanyaID INT NOT NULL,
  device_id VARCHAR(128) NOT NULL,
  push_device_id VARCHAR(128) NULL,
  device_name VARCHAR(120) NULL,
  -- VARCHAR et non ENUM : le client envoie aussi macOS / Windows / Linux, et
  -- ajouter une plateforme ne doit pas demander un ALTER TABLE.
  platform VARCHAR(20) NOT NULL DEFAULT 'unknown',
  ip_address VARCHAR(64) NULL,
  login_method ENUM('password','register','qr') NOT NULL,
  created_at DATETIME NOT NULL DEFAULT CURRENT_TIMESTAMP,
  last_active_at DATETIME NOT NULL DEFAULT CURRENT_TIMESTAMP,
  revoked_at DATETIME NULL,
  UNIQUE KEY uq_appareils_user_device (alanyaID, device_id),
  KEY idx_appareils_user (alanyaID, revoked_at),
  CONSTRAINT fk_appareils_user FOREIGN KEY (alanyaID)
    REFERENCES users(alanyaID) ON DELETE CASCADE
) ENGINE=InnoDB DEFAULT CHARSET=utf8mb4;
\end{lstlisting}

\begin{alerte}[Ordre de déploiement --- la migration d'abord, le code ensuite]
Les trois gardes d'authentification (\code{src/middleware/auth.js},
\code{src/middleware/authCustom.js}, \code{src/socket/handlers/auth.js}) joignent
\code{appareils} à \textbf{chaque} requête authentifiée. Déployer le code sans la table,
c'est un \code{ER\_NO\_SUCH\_TABLE} sur \emph{toute} l'API --- messages, appels et
\code{/auth/me} compris, pas seulement le QR. Le dépôt n'a pas de \emph{runner} de
migrations~: elles s'appliquent à la main, c'est donc à vérifier explicitement au
déploiement. L'avertissement figure en tête du fichier SQL.
\end{alerte}

\begin{note}[Pourquoi ni \code{user\_push\_devices} ni \code{userAccess}]
Deux autres tables évoquent déjà « appareil », chacune pour une raison différente.
\code{user\_push\_devices} (migration 020) a une portée strictement push --- heartbeat,
\code{appState}, jeton FCM~: des écritures fréquentes et peu sensibles. \code{userAccess}
est un journal en écriture seule, jamais muté, consultable seulement par un admin. Aucune
des deux n'est un état de session vivant et révocable. Mélanger \code{revoked\_at} --- un
champ sensible, rarement écrit --- avec le flux à haute fréquence de
\code{user\_push\_devices} aurait été le genre de raccourci qui coûte cher plus tard.
\end{note}

\begin{note}[\code{platform} en \code{VARCHAR(20)}, \code{login\_method} en \code{ENUM}]
Les deux colonnes sont un vocabulaire fermé~; une seule est déclarée comme tel, et
l'écart est délibéré. \code{platform} est \textbf{ouvert dans le temps}~: le client envoie
déjà \code{android}, \code{ios}, \code{web}, \code{macos}, \code{windows}, \code{linux} et
\code{unknown}, et une plateforme de plus ne doit pas coûter un \code{ALTER TABLE} sur une
table de session --- d'où le \code{VARCHAR(20)}, sur le modèle de
\code{user\_push\_devices.platform}. \code{login\_method}, lui, est \textbf{fermé par
nature}~: on ne se connecte que par mot de passe, par inscription ou par QR, et un
quatrième moyen serait de toute façon un chantier, pas une valeur de plus. La validation
n'est donc pas laissée à la seule base~: \code{deviceSessionService} normalise
\code{platform} contre une liste blanche (\code{\_normalizePlatform}, repli sur
\code{'unknown'}) et \textbf{lève} sur un \code{login\_method} inconnu, avant l'insertion.
\end{note}

\begin{alerte}[Pas de table pour la session éphémère]
La session de connexion (le pont entre le nouvel appareil et celui qui l'autorise) ne vit
\textbf{pas} dans une table SQL --- voir la justification au
chapitre~\ref{ch:qr-ephemere}, une fois son mécanisme introduit.
\end{alerte}

\chapter{Rendre la révocation réelle}
\label{ch:revocation}

\begin{constat}[Le point de départ, avant ce volet]
L'authentification était \textbf{100~\% sans état}~: un jeton d'accès (15~min) et un
jeton de rafraîchissement (30~jours, secret distinct), tous deux de simples JWT signés.
Aucune table \code{sessions} ni \code{tokens} n'existait, et
\code{POST~/api/auth/\allowbreak refresh} (\code{authCustomController.js}) se contentait de
vérifier la signature pour réémettre une nouvelle paire --- sans jamais rien invalider. Un
jeton de rafraîchissement qui fuyait restait valide 30~jours, quoi qu'il arrive. Ce
chapitre est la réponse à ce constat~: la même route refuse désormais de renouveler quoi
que ce soit si l'\code{appareilId} du jeton désigne une ligne révoquée.
\end{constat}

\begin{constat}[Chiffré]
Sur 78~comptes en production, \textbf{69} ont encore \code{users.device\_ID = 'INDEFINI'}
(la valeur par défaut de la colonne) et seulement \textbf{9} portent un vrai identifiant ---
hérité non pas du login, mais d'un effet de bord de
\code{pushDevicesController.js::registerPushDevice}, qui écrase par rétrocompatibilité
l'ancienne colonne mono-appareil à chaque enregistrement push réussi. Ce chemin ne couvre
ni les comptes qui refusent les notifications, ni le moment précis de la connexion. La
table \code{appareils} ne peut donc pas être un sous-produit d'un autre flux~: elle doit
être alimentée \textbf{directement} par \code{register}, \code{login} et l'approbation
d'une session QR.
\end{constat}

\begin{alerte}[Table dédiée, pas une extension de \code{user\_push\_devices}]
Réutiliser \code{user\_push\_devices} pour porter aussi la session d'authentification a été
envisagé, puis écarté~: les deux concepts ont des cycles de vie incompatibles sur une même
ligne. Un enregistrement push doit pouvoir survivre à la révocation d'une session --- c'est
précisément ce qui permet de notifier un appareil qu'il vient d'être déconnecté. Les
fusionner aurait rendu cette notification impossible au moment même où elle est la plus
utile.
\end{alerte}

Le JWT est étendu d'un champ \code{appareilId}~:

\begin{lstlisting}[language=JavaScript, caption={\code{src/middleware/authCustom.js} --- extension du payload}]
const generateAccessToken = (payload) =>
  jwt.sign({ ...payload, type: 'access' }, JWT_SECRET, { expiresIn: '15m' });
  // payload devient { alanyaID, email, appareilId }
\end{lstlisting}

La vérification de \code{appareils.revoked\_at IS NULL} se fait dans les gardes
d'authentification, à \textbf{chaque} requête --- pas seulement au rafraîchissement. Elles
sont \textbf{trois}, toutes sous \code{src/}~: \code{middleware/\allowbreak
authCustom.js}, \code{middleware/\allowbreak auth.js} et \code{socket/\allowbreak
handlers/\allowbreak auth.js}. Oublier la troisième laisserait un appareil révoqué en
401 sur tout le REST mais parfaitement vivant en temps réel~: il continuerait de recevoir
les messages du compte.

\begin{lstlisting}[language=SQL, caption={La jointure commune aux trois gardes}]
SELECT u.alanyaID, u.alanyaPhone, u.email
  FROM users u
  LEFT JOIN appareils a ON a.id = ? AND a.alanyaID = u.alanyaID
 WHERE u.alanyaID = ? AND u.exclus = 0
   AND (a.id IS NULL OR a.revoked_at IS NULL);
\end{lstlisting}

\begin{note}[Pourquoi une jointure \emph{externe}, et ce qu'elle concède]
Le \code{LEFT JOIN} et la clause \code{a.id IS NULL} laissent passer les jetons émis
\textbf{avant} la migration~026, qui ne portent pas d'\code{appareilId} --- sans quoi le
déploiement déconnecterait tout le monde d'un coup. La contrepartie est explicite~: la
révocation n'est réelle que pour les sessions ouvertes \emph{après} la migration. Les
anciennes s'éteignent d'elles-mêmes à l'expiration de leur jeton de rafraîchissement, au
plus tard 30~jours après.
\end{note}

\begin{note}[Le coût assumé]
Cela ajoute un \code{SELECT} indexé (\code{idx\_appareils\_user}) sur le chemin le plus
chaud de l'application --- en pratique, une jointure de plus sur une requête déjà faite.
Alternative plus légère~: ne vérifier qu'au rafraîchissement, avec un délai de grâce allant
jusqu'à 15~minutes. Choix retenu~: vérifier à chaque requête. Un bouton « Déconnecter » qui
met jusqu'à un quart d'heure à faire effet se comporterait, du point de vue de
l'utilisateur, comme un bouton cassé.
\end{note}

\section{Ce que fait vraiment la révocation}

\code{DELETE~/api/auth/device-sessions/:id} pose \code{revoked\_at = NOW()}, puis enchaîne
deux gestes de nature très différente~:

\begin{enumerate}[leftmargin=1.6em]
  \item \textbf{Un signal de confort} --- l'événement \code{auth:\allowbreak
        device\_revoked} part vers \emph{tout} le compte, dans la room
        \code{user\_\allowbreak \$\{alanyaID\}}, avec l'\textbf{identifiant matériel} de
        l'appareil coupé. Chaque client compare cette
        valeur à son propre \code{currentHardwareId()}~: seul l'appareil visé se
        reconnaît, et il efface alors sa session locale --- stockage sécurisé compris ---
        au lieu d'attendre son prochain 401. Diffuser à tout le compte plutôt que cibler
        une socket est
        volontaire~: l'appareil révoqué peut très bien ne pas avoir de socket ouverte à cet
        instant, et son identifiant matériel ne dit rien d'exploitable aux autres appareils
        du \emph{même} compte.
  \item \textbf{La garantie dure} --- le 401 des trois gardes. Elle ne dépend d'aucune
        coopération du client. \code{disconnectAppareilSockets}
        (\code{src/utils/userSocketRegistry.js}) ferme en outre les sockets encore ouvertes
        de cet appareil, \emph{après} l'\code{emit} pour que l'événement parte avant la
        fermeture.
\end{enumerate}

\begin{alerte}[La clé de fermeture des sockets est \code{appareilId}, pas \code{device\_id}]
\code{disconnectAppareilSockets} filtre sur \code{socket.appareilId} --- l'identifiant de
\emph{ligne}, porté par le JWT et posé sur la socket à son authentification. Filtrer sur un
\code{socket.deviceId} ne marcherait pas~: cette propriété-là porte l'UUID
\textbf{applicatif} du push, qui n'a aucun rapport avec \code{appareils.device\_id}
(section~\ref{sec:qr-device-id}). Les deux champs se ressemblent assez pour qu'on les
confonde, et l'erreur serait silencieuse~: aucune socket ne correspondrait, aucune
exception ne serait levée, et la fermeture ne ferait simplement rien.
\end{alerte}

\begin{alerte}[Un appareil révoqué ne se \emph{réactive} jamais]
Quand un appareil précédemment révoqué se reconnecte, \code{deviceSessionService.\allowbreak
recordLogin} ne remet pas \code{revoked\_at} à \code{NULL}~: il \textbf{supprime} la ligne
morte et en insère une neuve. La raison tient à l'\code{id}. Réactiver la ligne
conserverait son identifiant, donc l'\code{appareilId} porté par les anciens jetons~: le
jeton de rafraîchissement de 30~jours qu'on venait précisément de couper redeviendrait
valide, et l'appareil volé retrouverait l'accès sans rien faire d'autre qu'attendre une
reconnexion légitime du propriétaire. Un \code{id} neuf ne figure dans aucun jeton déjà
émis.
\end{alerte}

\begin{alerte}[Pas de ligne \code{appareils}, pas de jeton]
\code{register} et \code{login} exigent un identifiant d'appareil
(\code{hardware\_id}, à défaut \code{device\_ID}) et refusent la valeur sentinelle
\code{'INDEFINI'}~: sans lui, c'est un 400. Et si \code{recordLogin} échoue malgré tout,
ils répondent \textbf{503} sans émettre le moindre jeton, plutôt que de laisser passer une
connexion. Une session sans ligne \code{appareils} serait en effet invisible dans
« Appareils connectés » \emph{et} irrévocable --- exactement ce que ce chapitre entier
cherche à éliminer. Mieux vaut une connexion qui échoue franchement qu'une session
fantôme. Le même refus s'applique à l'approbation d'une session QR.
\end{alerte}

\chapter{Ajout de contact instantané}
\label{ch:qr-contact}

\code{addPreferredContact} (\code{src/controllers/preferredContactController.js}) faisait
déjà exactement ce qu'il faut~: rejet si l'utilisateur se scanne lui-même, vérification que
la cible existe et n'est pas exclue, vérification de non-blocage, vérification de
non-appartenance déjà acquise, puis \code{INSERT}. Cette logique est extraite en service
--- \code{src/services/contactService.js}, fonction \code{addContactByFriendId} ---
partagé par le contrôleur REST existant \textbf{et} par le résolveur QR, sans dupliquer les
règles. Le même fichier exporte \code{sanitizeUrl}, le filtre des sentinelles historiques
de \code{users.avatar\_url} (\code{'NON DEFINI'}, \code{'null'}\ldots{}), pour la même
raison~: recopié, il diverge et l'avatar s'affiche différemment d'un écran à l'autre.

\begin{note}[Le service ne connaît aucun code HTTP]
\code{addContactByFriendId} renvoie un résultat discriminé ---
\code{\{ok, reason, contact\}} avec \code{reason} $\in$ \code{added}, \code{self},
\code{not\_found}, \code{blocked}, \code{already} --- et laisse chaque appelant choisir sa
sémantique. C'est indispensable ici, parce que les deux appelants \emph{divergent
volontairement} sur deux cas~: là où la route REST \code{POST~/api/contacts/:id} répond
400 (soi-même) et 409 (déjà contact), la résolution d'un QR répond 200 dans les deux cas.
Un service qui aurait choisi les codes à leur place aurait forcé le QR à mentir sur ce qui
s'est passé.
\end{note}

\begin{center}
\begin{tabularx}{\linewidth}{@{}p{7.1cm}X@{}}
\toprule
\thd{Réponse de \mPOST{}~\code{/api/qr/resolve}} & \thd{Cas} \\
\midrule
\code{\{type:'contact', added:true,} \code{alreadyContact:false, contact\}} & ajouté \\
\code{\{type:'contact', added:false,} \code{alreadyContact:true, contact\}} & déjà favori
--- succès doux, aucune écriture \\
\code{200 \{type:'contact', self:true\}} & code correspondant à soi-même \\
\code{403} & bloqué, dans un sens ou dans l'autre \\
\code{404} & jeton inconnu, ou compte exclu \\
\bottomrule
\end{tabularx}
\end{center}

\begin{alerte}[Se scanner soi-même est un 200, pas un 400]
La partie~I range ce cas parmi les \textbf{succès doux}~: « rien ne se passe, un message
discret l'indique ». Une erreur~400 obligerait le client à distinguer, dans une même
branche d'échec, le code illisible (400 lui aussi) de son propre code --- et le moindre
raccourci d'affichage transformerait un geste anodin en message d'erreur rouge. Le
serveur répond donc 200 avec un drapeau \code{self}, que le client lit explicitement
(\code{isSelf} dans \code{lib/models/qr\_models.dart}). Même logique pour un contact déjà
enregistré~: le bandeau s'affiche, sans nouvelle écriture. Le champ s'appelle bien
\code{contact} et non \code{user}~: c'est la ligne \code{preferredContact} telle qu'elle
apparaîtra dans la liste, \code{idPrefContact} compris, pas la fiche de l'utilisateur.
\end{alerte}

\chapter{La session de connexion éphémère}
\label{ch:qr-ephemere}

\section{Une structure en mémoire, pas une table}

\begin{constat}[Vérifié dans le déploiement]
\code{Alanya-Backend} tourne en \textbf{un seul processus Node} --- \code{npm start} lance
\code{node server.js}, le déploiement fait \code{pm2 restart alanya-backend} sans mode
cluster, et aucun \code{ecosystem.config.js} n'existe. Un état partagé uniquement en
mémoire est donc visible de toutes les requêtes, sans risque qu'une requête tombe sur un
autre processus qui ignore la session.
\end{constat}

La session de connexion vit dans une \code{Map} en mémoire, sur le modèle exact de
\code{src/\allowbreak socket/\allowbreak state/\allowbreak pendingCalls.js}~:
\code{TTL\_MS}, \code{createdAt}~/ \code{expiresAt}, et un nettoyage porté par la structure
elle-même plutôt que par un job externe. Nouveau fichier~: \code{src/\allowbreak
socket/\allowbreak state/\allowbreak qrLoginSessions.js}.

\begin{note}[Ce que ça évite]
Ni migration, ni colonne, ni politique de rétention à trancher (une ligne de table qui
traînerait des semaines n'existe plus, par construction --- l'entrée disparaît de la
\code{Map} à l'expiration). Ce que ça coûte, en échange~: une session en cours ne survit
pas à un redémarrage du serveur --- l'utilisateur relance simplement le QR, un cas déjà
géré par la régénération automatique ci-dessous.
\end{note}

\begin{note}[Une \code{Map} bornée, pas seulement expirante]
L'expiration est \emph{paresseuse}~: pas de \code{setInterval}, une entrée périmée
disparaît à la première lecture qui la rencontre. Cela suffit tant qu'on lit ; or la route
de création n'est pas authentifiée, et son \code{authLimiter} a
\code{skipSuccessfulRequests}~: les créations \emph{réussies} ne consomment aucun quota.
Un client qui créerait sans jamais interroger ferait donc croître la \code{Map}
indéfiniment, dans le processus qui héberge aussi tous les sockets. Chaque création balaie
donc un échantillon d'entrées (\code{SWEEP\_PER\_CREATE}) et évince en FIFO au-delà de
\code{MAX\_SESSIONS} --- une \code{Map} JavaScript itère dans l'ordre d'insertion, l'éviction
est donc gratuite.
\end{note}

\section{Deux secrets par session, jamais interchangeables}
\label{sec:qr-deux-secrets}

C'est la pièce la moins évidente du dispositif, et celle dont dépend tout le reste~:
\textbf{le QR ne suffit pas à récupérer les jetons}.

\begin{center}
\begin{tabularx}{\linewidth}{@{}p{2.5cm}p{3.1cm}X@{}}
\toprule
\thd{Secret} & \thd{Dans le QR~?} & \thd{À qui, et pour quoi} \\
\midrule
\code{scanSecret} & \textbf{oui} & À l'appareil qui \textbf{scanne}~: il prouve qu'on a vu
le code. Exigé par \code{/resolve}, \code{/approve} et \code{/deny}. \\
\code{pollToken} & \textbf{jamais} & À l'appareil qui a \textbf{créé} la session, et à lui
seul --- remis une unique fois, dans la réponse à la création. Exigé pour interroger le
statut et retirer les jetons. \\
\bottomrule
\end{tabularx}
\end{center}

\begin{alerte}[Sans cette séparation, photographier le QR suffirait]
Avec un secret unique, quiconque prend en photo l'écran du nouvel appareil pourrait
interroger la session en boucle et \textbf{récupérer les jetons à la place du demandeur},
à l'instant précis où l'utilisateur approuve --- l'écran de confirmation, lui, aurait
montré un appareil parfaitement légitime. Le \code{pollToken} coupe cette attaque à la
racine~: il n'a jamais été affiché nulle part.
\end{alerte}

\begin{note}[En-tête \code{X-Poll-Token}, et non paramètre d'URL]
\code{GET~/api/auth/qr-session/:sessionId} lit le \code{pollToken} dans un \textbf{en-tête}
HTTP. Une \emph{query string} serait fonctionnellement identique et discrètement
désastreuse~: nginx et Cloudflare journalisent l'URL complète par défaut, et la réponse à
cette requête précise transporte \code{accessToken} et \code{refreshToken}. Le
\code{sessionId} étant public --- il est dans le QR --- le \code{pollToken} est le
\emph{seul} facteur~: il ne doit pas finir en clair dans un journal d'accès.
\end{note}

\begin{note}[Livraison des jetons à usage unique]
Dès que l'interrogation renvoie \code{approved} avec la paire de jetons, l'entrée est
\textbf{effacée} de la \code{Map} (\code{clear}). Les jetons ne repassent jamais par ce
canal~: un second appel, avec le même \code{pollToken}, ne trouve plus rien. Une session
inconnue et un \code{pollToken} faux donnent d'ailleurs la même réponse
--- \code{\{status:'expired'\}} --- pour qu'on ne puisse pas, avec le seul
\code{sessionId} public, sonder l'existence d'une session vivante. Même principe côté
scan~: un \code{scanSecret} invalide et une session inconnue renvoient tous deux 404.
\end{note}

\section{Cycle de vie}

\begin{center}
\begin{tabularx}{\linewidth}{@{}p{2.4cm}X@{}}
\toprule
\thd{Statut} & \thd{Déclencheur} \\
\midrule
\code{pending}   & créée par le nouvel appareil, \code{expiresAt} $=$ maintenant $+$ 90~s \\
\code{scanned}   & un appareil authentifié a résolu le jeton \\
\code{approving} & réservation interne pendant l'approbation, jamais exposée \\
\code{approved}  & confirmation explicite --- crée la ligne \code{appareils}, puis livre
les jetons une seule fois \\
\code{denied}    & refus explicite. L'entrée est \textbf{conservée} jusqu'à son TTL~: la
supprimer afficherait « expiré » au demandeur, un message trompeur \\
\code{expired}   & statut \emph{synthétique} --- l'entrée n'est plus dans la \code{Map},
qu'elle ait expiré ou déjà livré ses jetons \\
\bottomrule
\end{tabularx}
\end{center}

À expiration, le nouvel appareil régénère seul une nouvelle session --- l'utilisateur ne
voit jamais un code mort.

\begin{alerte}[Réservation atomique de l'approbation]
L'approbation n'est pas instantanée~: elle enchaîne plusieurs \code{await} --- lecture du
compte, écriture dans \code{appareils}, signature des jetons. Deux confirmations quasi
simultanées franchiraient donc toutes deux la garde de statut, et la \textbf{seconde
écraserait le résultat de la première}~: l'appareil demandeur ouvrirait le compte du
second approbateur. \code{qrLoginSessions.beginApproval()} bascule le statut en
\code{approving} \emph{avant} le premier \code{await}~; comme rien n'est asynchrone entre
la lecture et l'écriture du statut, la transition est atomique en JavaScript mono-thread.
La seconde confirmation reçoit alors un 409. \code{abortApproval()} rend la session à son
statut précédent si l'approbation échoue en cours de route --- compte introuvable,
\code{recordLogin} en échec --- pour qu'un incident ne condamne pas un code encore valide.
\end{alerte}

\begin{note}[Le cas où la session expire \emph{pendant} l'approbation]
\code{approve()} peut échouer après que la ligne \code{appareils} a été créée. Le
contrôleur teste donc son retour~: en cas d'échec, il \textbf{supprime la ligne} qu'il
vient d'insérer et répond 410 \code{QR\_SESSION\_EXPIRED}. Sans ce test, on annoncerait
« approuvé », on émettrait une paire de jetons que personne ne viendrait chercher, et on
laisserait un appareil orphelin dans la liste « Appareils connectés ».
\end{note}

\section{L'exception d'authentification socket}
\label{sec:qr-socket-exception}

\begin{constat}
Chaque handler socket existant --- chat, appels, réunions, accusés de réception --- ouvre
par la même garde~: \code{if~(!socket.authenticated)~return;}. C'est la règle universelle
du fichier \code{server.js} et de \code{src/socket/handlers/}.
\end{constat}

Le nouvel appareil n'est, par définition, pas encore authentifié quand il attend une
approbation. Un nouveau handler \code{src/\allowbreak socket/\allowbreak
handlers/\allowbreak qrLogin.js} déroge donc --- seule exception du dépôt --- à cette
garde~: \code{socket.on(\allowbreak 'qr\_session:join', \ldots{})} valide directement
contre la \code{Map} de \code{qrLoginSessions.js} plutôt que contre une session
utilisateur, puis rejoint la room \code{qr\_session\_\allowbreak \$\{sessionId\}}.

\begin{note}[Une exception bornée par construction]
Elle exige le \textbf{\code{pollToken}} (section~\ref{sec:qr-deux-secrets}), pas le
\code{scanSecret}~: photographier le QR ne suffit donc pas à entrer dans la room. Le socket
ne devient jamais \code{authenticated}, ne rejoint jamais \code{user\_*} et ne reçoit donc
aucun autre événement du serveur. La room ne transporte que le \emph{statut}, jamais de
jeton --- la livraison des jetons reste sur le \emph{poll} HTTP à usage unique. Et la
session s'auto-détruit au bout de 90~secondes. Quatre bornes, aucune reposant sur la
bonne volonté du client.
\end{note}

\begin{note}[Le repli FCM ne fonctionne pas ici]
\code{notificationService.js::sendToUser} cible par \code{alanyaID} --- une donnée que le
nouvel appareil n'a, par construction, pas encore. Le repli habituel (push en complément du
socket, comme pour les appels) est donc \textbf{impossible} sur ce chemin précis. Le seul
canal qui fonctionne avant authentification est celui que l'appareil a lui-même ouvert~: la
room socket, ou le polling REST.
\end{note}

\begin{alerte}[Décision~v1~: polling seul côté mobile]
Le socket authentifié de \code{TalkyApiClient} suppose systématiquement un jeton d'accès
valide --- toute la mécanique de reconnexion et de \emph{watchdog} est bâtie autour d'une
session déjà établie. Faire cohabiter un mode « anonyme, scellé à un jeton de 90~s » dans
cette même classe est un risque de régression sur un composant déjà complexe. \textbf{Le
client mobile v1 n'utilise que le polling REST} (\code{GET~/api/auth/qr-session/:sessionId},
toutes les 2~s) --- le délai perçu est négligeable face au temps humain de déverrouiller
l'autre téléphone et de confirmer. La room \code{qr\_session\_\$\{sessionId\}} reste
implémentée côté serveur~: elle ne coûte rien et prépare la voie web
(chapitre~\ref{ch:qr-web}), qui peut se permettre un socket ouvert en continu.
\end{alerte}

\begin{note}[\code{ttlMs} plutôt qu'une simple date de fin]
La création renvoie \code{expiresAt} \textbf{et} \code{ttlMs}. Le client décompte à partir
de \code{ttlMs} et de l'instant de \emph{réception}, sans jamais comparer son horloge à
celle du serveur. Un téléphone dont l'horloge avance de plus de 90~secondes --- ce qui
n'a rien d'exotique --- verrait sinon son code expiré dès son affichage, et le
régénérerait en boucle sans jamais laisser le temps de le scanner.
\end{note}

\begin{alerte}[Confirmation explicite, jamais d'approbation automatique]
Rappel technique du choix produit~: \code{POST~/api/qr/resolve} sur un jeton de type
\code{login} ne fait que passer le statut à \code{scanned} et renvoyer les informations de
l'appareil demandeur --- il est \emph{strictement en lecture seule} côté autorisation. Seul
\code{POST~/api/auth/qr-session/:sessionId/approve}, appelé après une action explicite de
l'utilisateur sur un écran dédié, crée la ligne \code{appareils} et émet les jetons. Sans
cette séparation en deux appels, le scan seul suffirait à connecter --- exactement le
scénario de « quishing » que la partie~I met en garde.
\end{alerte}

\begin{alerte}[Les informations affichées viennent d'un appelant non authentifié]
\code{deviceName} et \code{platform} sont fournis par l'appareil qui \emph{crée} la
session --- lequel n'est, par construction, pas encore connecté. Ils s'affichent pourtant
tels quels sur l'écran de confirmation, seul rempart contre le quishing. La plateforme est
donc contrainte à un vocabulaire fermé \textbf{dès la création} (et non au moment de
l'insertion en base, trop tard car après l'approbation), et le libellé libre est borné à
60~caractères puis présenté côté client comme « déclaré par l'appareil », jamais comme un
fait vérifié.
\end{alerte}

\chapter{API}

Deux routeurs seulement~: \code{src/\allowbreak routes/\allowbreak qr.js}, monté sur
\code{/api/qr} par \code{server.js}, et \code{src/\allowbreak routes/\allowbreak
authCustom.js}, qui accueille les chemins \code{/api/\allowbreak auth/\allowbreak
qr-session*} et \code{/api/\allowbreak auth/\allowbreak device-sessions*} --- exactement
comme il accueille déjà \code{push-devices} et \code{notification-prefs}.

\begin{center}
\begin{tabularx}{\linewidth}{@{}>{\raggedright\arraybackslash}p{5.9cm}
                             >{\raggedright\arraybackslash}p{3.2cm}X@{}}
\toprule
\thd{Endpoint} & \thd{Garde} & \thd{Effet} \\
\midrule
\mGET{}~\code{/api/qr/me} & \code{authCustom} &
\emph{mint-if-absent} de \code{qr\_public\_id} \\
\mPOST{}~\code{/api/qr/me/\allowbreak regenerate} & \code{authCustom} &
nouveau jeton, l'ancien devient invalide \\
\mPOST{}~\code{/api/qr/resolve} & \code{authCustom} $+$ \code{qrResolveLimiter} &
résout un jeton scanné (contact ou login) \\
\mPOST{}~\code{/api/auth/\allowbreak qr-session} & \code{authLimiter} (public) &
crée une session~; renvoie \code{scanSecret}, \code{pollToken}, \code{payload},
\code{expiresAt}, \code{ttlMs} \\
\mGET{}~\code{/api/auth/\allowbreak qr-session/:sessionId} & en-tête
\code{X-Poll-Token} (public) & interrogation par le nouvel appareil~; livre les jetons
une seule fois \\
\mPOST{}~\code{/api/auth/\allowbreak qr-session/:sessionId/\allowbreak approve} &
\code{authCustom} $+$ \code{scanSecret} & crée \code{appareils}, émet les jetons \\
\mPOST{}~\code{/api/auth/\allowbreak qr-session/:sessionId/\allowbreak deny} &
\code{authCustom} $+$ \code{scanSecret} & marque \code{denied} \\
\mGET{}~\code{/api/auth/\allowbreak device-sessions} & \code{authCustom} &
liste des appareils non révoqués \\
\mDEL{}~\code{/api/auth/\allowbreak device-sessions/:id} & \code{authCustom} &
révocation immédiate \\
\bottomrule
\end{tabularx}
\end{center}

\begin{center}
\begin{tabularx}{\linewidth}{@{}p{6.6cm}X@{}}
\toprule
\thd{Module} & \thd{Rôle} \\
\midrule
\code{src/controllers/qrController.js} & \code{/api/qr/*} --- identité et résolution \\
\code{src/controllers/qrAuthController.js} & \code{/api/\allowbreak auth/\allowbreak
qr-session*} \textbf{et} \code{/api/\allowbreak auth/\allowbreak device-sessions*}~: un
seul contrôleur, parce que les deux familles manipulent le même objet --- la session d'un
appareil --- l'une pour l'ouvrir, l'autre pour la fermer \\
\code{src/services/deviceSessionService.js} & \code{recordLogin} /
\code{touchLastActive} / \code{normalizePlatform} --- la table \code{appareils}, partagée
avec \code{login} et \code{register} \\
\code{src/services/contactService.js} & \code{addContactByFriendId}, partagé avec
\mPOST{}~\code{/api/\allowbreak contacts/:id} \\
\code{src/utils/qrToken.js} & jetons opaques, payloads, comparaison à temps constant \\
\code{src/socket/state/qrLoginSessions.js} & la \code{Map} des sessions éphémères \\
\bottomrule
\end{tabularx}
\end{center}

\section{Limitation de débit}

\begin{center}
\begin{tabularx}{\linewidth}{@{}>{\raggedright\arraybackslash}p{4.9cm}
                             >{\raggedright\arraybackslash}p{3.4cm}X@{}}
\toprule
\thd{Route} & \thd{Limiteur} & \thd{Pourquoi ce calibre} \\
\midrule
\mPOST{}~\code{/api/auth/\allowbreak qr-session} & \code{authLimiter}
(10~/~15~min~/~IP) & Création, \emph{une} par code affiché. Le limiteur a
\code{skipSuccessfulRequests}~: seuls les échecs consomment le quota. \\
\mPOST{}~\code{/api/qr/resolve} & \code{qrResolveLimiter} (30~/~min~/~IP) &
Un humain qui scanne dépasse rarement quelques codes par minute. \\
\mGET{}~\code{/api/auth/\allowbreak qr-session/:sessionId} & un limiteur
\textbf{dédié}, de l'ordre de 60~/~min~/~IP & Interrogation toutes les 2~s pendant
90~s~: environ 45~appels \emph{par code affiché}. \\
\bottomrule
\end{tabularx}
\end{center}

\begin{alerte}[Ne jamais poser \code{authLimiter} sur l'interrogation de statut]
Le réflexe serait d'aligner l'interrogation sur la création, toutes deux publiques. Il
casserait la fonctionnalité~: \code{authLimiter} autorise 10~requêtes par 15~minutes,
quand un seul code en consomme environ 45 en 90~secondes. L'utilisateur verrait son écran
se figer au bout de vingt secondes, puis échouer --- et le symptôme ressemblerait à une
panne réseau, pas à un plafond. Cette route a besoin d'un limiteur \emph{dimensionné pour
du polling}, de l'ordre de 60~requêtes par minute et par~IP~: assez large pour deux
appareils derrière un même NAT, assez étroit pour rendre coûteux le sondage massif de
\code{sessionId}. \textbf{À ce jour, aucun limiteur dédié n'est posé sur cette route}~: elle
n'est couverte que par le \code{generalLimiter} global (300~/~min~/~IP), qui la protège
d'un abus grossier mais n'est pas calibré pour elle.
\end{alerte}

\chapter{Client Flutter}

\section{Couche données}

\begin{center}
\begin{tabularx}{\linewidth}{@{}p{6.4cm}X@{}}
\toprule
\thd{Fichier} & \thd{Contenu} \\
\midrule
\code{lib/api/qr\_api.dart} \emph{(nouveau)} & Extension \code{QrApi on TalkyApiClient},
\code{part of} le client, sur le modèle de \code{auth\_api.dart} / \code{users\_api.dart}~:
\code{getMyQr}, \code{regenerateQr}, \code{resolveQr}. \\
\code{lib/api/qr\_auth\_api.dart} \emph{(nouveau)} & Extension \code{QrAuthApi}~:
\code{createQrSession}, \code{pollQrSession} (qui pose l'en-tête \code{X-Poll-Token}),
\code{approveQrSession}, \code{denyQrSession}, \code{listDeviceSessions},
\code{revokeDeviceSession}. Séparée de la précédente parce qu'elle vise
\code{/api/auth/*} et qu'une partie de ses appels se fait \textbf{sans} jeton d'accès. \\
\code{lib/models/qr\_models.dart} \emph{(nouveau)} & \code{QrIdentity},
\code{QrResolveResult} (\code{isSelf}, \code{added}, \code{alreadyContact}),
\code{QrLoginCreated}, \code{QrLoginPollResult}, \code{AppareilSession}. \\
\code{lib/providers/auth\_provider.dart} \emph{(corrigé)} & \code{login()} /
\code{register()} transmettent désormais \code{hardwareId:
await~TalkyApiClient.\allowbreak currentHardwareId()} --- l'identifiant matériel du
chapitre~\ref{ch:qr-donnees}, jusqu'ici absent malgré le paramètre déjà accepté. \\
\code{lib/api/socket\_api.dart} \emph{(modifié)} & Écoute \code{auth:device\_revoked} et
\code{auth:conflict}~; sur le premier, compare l'identifiant reçu à son propre
\code{currentHardwareId()} et se déconnecte s'il se reconnaît. \\
\bottomrule
\end{tabularx}
\end{center}

\begin{note}[Aucune table Drift pour ce volet]
Ni le QR personnel ni la liste d'appareils ne passent par le cache hors ligne. Le premier
est bon marché à retélécharger et sans utilité hors ligne --- scanner exige déjà un réseau.
La seconde est une donnée de \textbf{sécurité}~: afficher un appareil déjà révoqué, ou
masquer un appareil suspect récent, à cause d'un cache périmé, serait activement trompeur.
Pas de nouvelle table dans \code{lib/core/db/app\_database.dart}, pas de montée de
\code{schemaVersion}.
\end{note}

\section{Interface}

\begin{center}
\begin{tabularx}{\linewidth}{@{}p{8.4cm}X@{}}
\toprule
\thd{Fichier} & \thd{Objet} \\
\midrule
\code{screens/profile/\allowbreak qr\_scanner\_screen.dart} \emph{(nouveau)} & Scanner
plein écran (\code{mobile\_scanner}), gabarit
\code{screens/chats/\allowbreak camera\_screen.dart}. Autonome~: appelle \code{resolveQr},
branche selon \code{type}. \\
\code{screens/profile/qr\_code\_screen.dart} \emph{(nouveau)} & « Mon code »~: deux
onglets, \emph{Mon code} et \emph{Scanner}, partage et régénération. \\
\code{widgets/alanya\_qr\_view.dart} \emph{(nouveau)} & Rendu unique des QR --- identité
comme connexion. \\
\code{screens/profile/\allowbreak qr\_login\_confirm\_screen.dart} \emph{(nouveau)} & Fiche
de l'appareil scanné, Confirmer / Refuser. \\
\code{screens/authentification/\allowbreak qr\_login\_screen.dart} \emph{(nouveau)} & QR de
session, compte à rebours fondé sur \code{ttlMs}, \code{Timer.periodic} de \emph{polling}
toutes les 2~s. \\
\code{screens/profile/\allowbreak connected\_devices\_screen.dart} \emph{(nouveau)} & Liste
\code{appareils}. Le repère « Cet appareil » vient du drapeau \code{current} \textbf{calculé
par le serveur} (comparaison de \code{appareils.id} à l'\code{appareilId} du jeton), pas
d'une comparaison locale. \\
\code{screens/profile/profile\_screen.dart} & Icône QR dans l'en-tête. \\
\code{widgets/add\_contact\_sheet.dart} & Bouton « Scanner un code ». \\
\code{screens/profile/\allowbreak account\_security\_screen.dart} & Sections « Appareils
connectés » et « Lier un appareil ». \\
\code{screens/authentification/login\_screen.dart} & Bouton « Se connecter avec un code
QR ». \\
\code{pubspec.yaml} & \code{+ qr\_flutter}, \code{+ mobile\_scanner}. \\
\code{lib/l10n/*} & Quatre familles de clés~: \code{qrMyCode*}, \code{qrScan*},
\code{qrLogin*} / \code{qrApprove*}, \code{qrDevices*}, plus \code{qrBannerNewDevice} et
\code{qrBannerSignedOutRemotely}. \\
\bottomrule
\end{tabularx}
\end{center}

\begin{note}[Rendu du QR]
\code{qr\_flutter} (\code{QrImageView}) permet \code{eyeStyle} et \code{dataModuleStyle}
personnalisés (points arrondis, couleur \code{brandPrimary}), un \code{embeddedImage} pour
le logo, et un niveau de correction d'erreur \code{QrErrorCorrectLevel.H} (30~\% de
redondance) --- nécessaire pour tolérer l'occlusion du logo central sans nuire à la
lecture. Le scan repose sur \code{mobile\_scanner} (CameraX et ML~Kit sur Android,
AVFoundation/Vision sur iOS), retenu plutôt que \code{qr\_code\_scanner} pour son socle
natif et son entretien plus actif dans l'écosystème Flutter récent.
\end{note}

\begin{note}[Un seul widget de rendu, et pourquoi]
\code{alanya\_qr\_view.dart} existe parce que les deux rendus avaient \emph{déjà} divergé
--- yeux ronds contre carrés, bicolore contre monochrome, deux fichiers de logo différents
--- donnant deux identités visuelles à la même fonctionnalité. Le fond blanc en thème
sombre et le niveau~H y sont posés une fois pour toutes~: ce ne sont pas des choix
esthétiques mais des conditions de lisibilité, et les laisser à chaque écran, c'était les
perdre au premier oubli.
\end{note}

\chapter{Vers une future version web}
\label{ch:qr-web}

Rien dans l'API n'est spécifique au mobile~: une session de connexion se crée, s'interroge
et s'approuve par de simples appels REST, et la room socket
\code{qr\_session\_\$\{sessionId\}} (section~\ref{sec:qr-socket-exception}) fonctionne
identiquement pour n'importe quel client. La colonne \code{appareils.platform} traite déjà
\code{'web'} comme une valeur de premier ordre --- exactement le même choix que fait
\code{user\_push\_devices.platform} depuis la migration~020.

Un onglet de navigateur serait même un \textbf{meilleur} client que l'application mobile
pour ce flux~: il reste ouvert en continu, ce qui rend le canal socket --- volontairement
laissé de côté côté mobile (section~\ref{sec:qr-socket-exception}) --- immédiatement utile
plutôt qu'à préparer pour plus tard. Seul le rendu du code changerait (un \code{canvas} ou
un \code{svg} plutôt que \code{qr\_flutter}). Aucun écran web n'est conçu dans ce volet.

\chapter{Impact fichier par fichier}

\section{Backend --- Alanya-Backend}

\begin{center}
\begin{tabularx}{\linewidth}{@{}p{9.0cm}p{1.3cm}X@{}}
\toprule
\thd{Fichier} & \thd{Nature} & \thd{Objet} \\
\midrule
\code{migrations/026\_qr\_appareils.sql} & nouveau & \code{qr\_public\_id},
\code{appareils} \\
\code{src/utils/qrToken.js} & nouveau & Jetons opaques, \code{QR\_BASE\_URL}, payloads,
\code{secretMatches} \\
\code{src/services/contactService.js} & nouveau & \code{addContactByFriendId} et
\code{sanitizeUrl}, extraits d'\code{addPreferredContact} \\
\code{src/services/deviceSessionService.js} & nouveau & \code{recordLogin},
\code{touchLastActive}, \code{normalizePlatform} \\
\code{src/socket/state/qrLoginSessions.js} & nouveau & \code{Map} en mémoire des sessions
de connexion, sur le modèle de \code{pendingCalls.js} \\
\code{src/controllers/qrController.js} & nouveau & \code{/api/qr/*} \\
\code{src/controllers/qrAuthController.js} & nouveau & \code{/api/auth/qr-session*} et
\code{/api/auth/device-sessions*} \\
\code{src/routes/qr.js} & nouveau & Monte \code{/api/qr} \\
\code{src/socket/handlers/qrLogin.js} & nouveau & \code{qr\_session:join}, seule exception
à la garde d'authentification \\
\code{server.js} & modifié & Montage des routes, enregistrement du handler socket \\
\code{src/routes/authCustom.js} & modifié & Routes \code{qr-session*} et
\code{device-sessions*} \\
\code{src/middleware/rateLimiter.js} & modifié & \code{qrResolveLimiter} \\
\code{src/middleware/authCustom.js} & modifié & Payload JWT étendu, vérification
\code{revoked\_at} \\
\code{src/middleware/auth.js} & modifié & Idem, pour rester cohérent avec
\code{authCustom.js} \\
\code{src/socket/handlers/auth.js} & modifié & Même vérification côté socket~;
\code{socket.appareilId} \\
\code{src/controllers/authCustomController.js} & modifié & \code{register} / \code{login}~:
\code{recordLogin} obligatoire, 503 sinon~; \code{refreshToken}~: refus si l'appareil est
révoqué \\
\code{src/controllers/preferredContactController.js} & modifié & Appelle le service
extrait, au lieu de porter la logique \\
\code{src/utils/userSocketRegistry.js} & modifié & \code{disconnectAppareilSockets},
filtrant sur \code{socket.appareilId} \\
\bottomrule
\end{tabularx}
\end{center}

\begin{note}[Trois fichiers que le dossier avait prévus et qui n'existent pas]
La conception initiale annonçait \code{src/constants/qrSessionStatus.js},
\code{src/services/qrIdentityService.js} et \code{src/services/qrLoginService.js}. Aucun
n'a été écrit, et c'est délibéré~: les statuts sont des chaînes déjà exhaustivement
énumérées par les transitions de \code{qrLoginSessions.js}, et les deux services
n'auraient fait que déléguer, l'un à \code{qrController}, l'autre à
\code{qrLoginSessions}. Trois indirections vides valent moins qu'un contrôleur qu'on lit
d'un bout à l'autre.
\end{note}

\section{Mobile --- Talky (Flutter)}

\begin{center}
\begin{tabularx}{\linewidth}{@{}p{9.0cm}p{1.3cm}X@{}}
\toprule
\thd{Fichier} & \thd{Nature} & \thd{Objet} \\
\midrule
\code{lib/api/qr\_api.dart} & nouveau & Identité et résolution \\
\code{lib/api/qr\_auth\_api.dart} & nouveau & Sessions de connexion et appareils \\
\code{lib/models/qr\_models.dart} & nouveau & Modèles de la couche données QR \\
\code{lib/screens/profile/qr\_code\_screen.dart} & nouveau & Mon code \\
\code{lib/screens/profile/qr\_scanner\_screen.dart} & nouveau & Scanner \\
\code{lib/widgets/alanya\_qr\_view.dart} & nouveau & Rendu unique des QR \\
\code{lib/screens/profile/qr\_login\_confirm\_screen.dart} & nouveau & Confirmation \\
\code{lib/screens/authentification/qr\_login\_screen.dart} & nouveau & QR de session \\
\code{lib/screens/profile/connected\_devices\_screen.dart} & nouveau & Mes appareils \\
\code{android/\ldots{}/alanya237/alanya/MainActivity.kt} & modifié & Canal natif
\code{com.alanya/device\_identity} $\rightarrow$ \code{ANDROID\_ID} \\
\code{lib/talky\_api\_client.dart} & modifié & \code{currentHardwareId()}, \code{part}
des deux nouvelles extensions \\
\code{lib/providers/auth\_provider.dart} & \textbf{corrigé} & \code{hardwareId} transmis au
login / register \\
\code{lib/api/auth\_api.dart} & modifié & Champ \code{hardware\_id} dans le corps \\
\code{lib/api/socket\_api.dart} & modifié & \code{auth:device\_revoked},
\code{auth:conflict} \\
\code{lib/talky\_models.dart} & modifié & \code{SocketEvents.\allowbreak
authDeviceRevoked}, \code{AccountDeviceLogin} \\
\code{lib/screens/profile/profile\_screen.dart} & modifié & Icône QR \\
\code{lib/widgets/add\_contact\_sheet.dart} & modifié & Bouton « Scanner un code » \\
\code{lib/screens/profile/account\_security\_screen.dart} & modifié & Sections « Appareils
connectés » et « Lier un appareil » \\
\code{lib/screens/authentification/login\_screen.dart} & modifié & Bouton « Se connecter
avec un code QR » \\
\code{pubspec.yaml} & modifié & \code{qr\_flutter}, \code{mobile\_scanner} \\
\code{lib/l10n/*} & modifié & Nouvelles clés \\
\bottomrule
\end{tabularx}
\end{center}

\begin{note}[Administration --- hors périmètre]
Aucun écran \code{talky-admin} n'est nécessaire à ce volet. Une vue support « appareils
connectés d'un compte » serait une extension naturelle mais n'a pas été demandée --- signalée
au chapitre suivant, pas conçue ici.
\end{note}

\chapter{Vérification}

\section{Backend}

Appliquer \code{026\_qr\_appareils.sql} \textbf{avant} de déployer le code, redémarrer,
puis~:

\begin{enumerate}[leftmargin=1.6em]
  \item \code{GET~/api/qr/me} deux fois $\rightarrow$ même jeton ; après
        \code{POST~/regenerate}, un jeton différent, et l'ancien renvoie 404 sur
        \code{resolve} ;
  \item scanner (résoudre) le jeton d'un autre compte $\rightarrow$ ligne
        \code{preferredContact} créée, 200 en le rejouant (succès doux, pas de doublon) ;
        scanner \emph{son propre} jeton $\rightarrow$ 200 \code{\{self:true\}}, jamais
        une erreur ;
  \item créer une session QR, la laisser expirer $\rightarrow$ \code{status = expired} au
        polling suivant ;
  \item interroger une session avec un \code{X-Poll-Token} erroné $\rightarrow$
        \code{expired}, indistinguable d'une session inconnue~; sans en-tête du tout
        $\rightarrow$ 400 ;
  \item créer une session, la résoudre depuis un second compte, \textbf{approuver}
        $\rightarrow$ ligne \code{appareils} créée, jetons valides, \code{appareilId}
        présent dans leur charge utile. Rejouer l'interrogation $\rightarrow$
        \code{expired}~: les jetons ne sont livrés qu'une fois ;
  \item approuver \textbf{deux fois} la même session, quasi simultanément $\rightarrow$
        la seconde reçoit 409, jamais une seconde paire de jetons ;
  \item \code{DELETE~/api/auth/device-sessions/:id} sur un appareil actif $\rightarrow$
        sa requête authentifiée \textbf{suivante} échoue immédiatement, socket fermée ;
        le même appareil qui se reconnecte obtient une ligne au \code{id} \textbf{différent}.
\end{enumerate}

\section{Client}

\code{flutter pub get} $\rightarrow$ \code{flutter run}.

\begin{enumerate}[leftmargin=1.6em]
  \item Scanner le code d'un contact $\rightarrow$ ajout instantané, bandeau « Annuler »
        fonctionnel pendant 5~s.
  \item Scanner son propre code, puis un code expiré $\rightarrow$ messages distincts,
        caméra jamais bloquée.
  \item Connecter un second appareil de bout en bout~: code affiché, scan depuis le
        premier, fiche correcte, confirmation, jetons reçus par polling.
  \item \textbf{Thème sombre}~: la carte du QR personnel reste blanche.
  \item Réinstaller l'app sur le même appareil physique~: sur \textbf{Android}, il ne
        réapparaît \emph{pas} comme un second appareil dans la liste (\code{ANDROID\_ID}
        stable). Sur iOS, il réapparaît si c'était la dernière application du vendeur ---
        limite connue.
  \item Déconnecter un appareil depuis un autre~: l'appareil visé se déconnecte
        \textbf{tout seul} à la réception d'\code{auth:device\_revoked}, sans attendre un
        appel réseau. Refaire le test \emph{application fermée}~: il doit alors tomber au
        premier 401.
  \item Régler l'horloge du téléphone avec 5~minutes d'avance, puis afficher un QR de
        connexion~: le compte à rebours doit partir de 90~s (\code{ttlMs}) et non expirer
        immédiatement.
\end{enumerate}

\chapter{Points de vigilance et questions ouvertes}

\begin{alerte}[Ajout instantané, sans garde-fou fort]
La règle produit --- ajouter sans confirmation --- n'a pour seul filet qu'un bandeau
annulable de 5~secondes. Elle est appliquée telle que demandée ; elle mérite d'être
révisée si les codes personnels se retrouvent affichés dans des contextes publics
(vitrines, affiches).
\end{alerte}

\begin{enumerate}[leftmargin=1.6em]
  \item \textbf{Coût de la vérification en garde} (chapitre~\ref{ch:revocation})~:
        une jointure de plus par requête authentifiée, dans \textbf{trois} fichiers quasi
        dupliqués (\code{auth.js}, \code{authCustom.js}, \code{socket/handlers/auth.js}).
        Assumé ici ; leur fusion serait une occasion naturelle, mais indépendante de ce
        volet --- et le risque, en l'état, est qu'une évolution future n'en corrige que
        deux sur trois.
  \item \textbf{Pas de limiteur dédié sur l'interrogation de statut.}
        \code{GET~/api/auth/qr-session/:sessionId} n'est couvert que par le
        \code{generalLimiter} global (300~/~min~/~IP). Un limiteur calibré pour du
        \emph{polling} (${\sim}60$~/~min~/~IP) reste à poser --- en aucun cas
        \code{authLimiter}, qui casserait la fonctionnalité.
  \item \textbf{\code{auth:conflict} a désormais un émetteur.} L'événement, qui existait
        côté mobile sans jamais être émis, est maintenant diffusé par
        \code{authCustomController.js} (connexion par mot de passe) et par
        \code{qrAuthController.js} (approbation d'une session QR). Son nom reste
        discutable~: il annonce une \emph{nouvelle connexion sur le compte}, pas un
        conflit. Le renommer reste ouvert, et coûterait un déploiement coordonné des deux
        côtés.
  \item \textbf{Aucune limite} au nombre d'appareils simultanés n'est proposée. Rien dans
        la demande ne l'impose ; un garde-fou futur (avertissement au-delà d'un seuil, à
        la manière de certains services bancaires) resterait compatible avec ce schéma.
  \item \textbf{\code{push\_device\_id} n'est écrit par personne.} La colonne existe,
        nullable, et \emph{aucun} code ne la renseigne à ce jour --- ni le client, ni le
        backend. Elle reste posée pour le badge « notifications activées », qui demanderait
        de faire remonter le \code{deviceId} push au moment du login. En attendant, l'écran
        « Appareils connectés » doit se comporter comme si elle était toujours vide (aucun
        badge, pas d'erreur), ce qui est le cas aujourd'hui. Rappel~: un utilisateur qui
        refuse les notifications n'aura de toute façon jamais de correspondance dans
        \code{user\_push\_devices}.
  \item \textbf{L'identifiant matériel n'est pas un secret.} \code{ANDROID\_ID} est
        diffusé à tout le compte dans \code{auth:device\_revoked}, et
        \code{identifierForVendor} est lisible par toute application du même vendeur. Cela
        reste acceptable~: connaître le \code{device\_id} d'un appareil ne permet ni de
        s'y substituer ni de le révoquer --- ces deux gestes demandent un jeton valide du
        compte. Mais c'est une raison de plus pour que l'API n'expose jamais ce champ dans
        la liste des appareils, ce qu'elle ne fait pas (pas plus que l'adresse~IP).
\end{enumerate}


\vfill
\begin{center}\small\color{AlanyaGray}\sffamily
Fin du dossier --- étape 6. \\
Volet indépendant --- peut se construire en parallèle des cinq premiers.
\end{center}

\end{document}
